\section*{Bab \ref{ch:g11:functional-groups} --- Gugus Fungsi dan Keluarga Senyawa}

\begin{solution}{exo:g11:functional-groups:1}
Hidroksil, \omterm{def:g11:functional-groups:alcohol}{alkohol}; karbonil di tengah rantai, \omterm{def:g11:functional-groups:carbonyl}{keton}; karboksil, \omterm{def:g11:functional-groups:carboxylic-acid}{asam
karboksilat}; gugus \omterm{def:g11:functional-groups:carboxylic-acid}{ester}, \omterm{def:g11:functional-groups:carboxylic-acid}{ester}; gugus amino \ce{-NH2}, \omterm{def:g11:functional-groups:amine}{amina}.
\end{solution}

\begin{solution}{exo:g11:functional-groups:2}
Metanol; propan-1-ol; asam metanoat; propanal; propanon.
\end{solution}

\begin{solution}{exo:g11:functional-groups:3}
\omterm{def:g11:functional-groups:alcohol}{Alkohol}; \omterm{def:g11:functional-groups:carbonyl}{aldehida}; \omterm{def:g11:functional-groups:carboxylic-acid}{ester}; \omterm{def:g11:functional-groups:amine}{amida}; \omterm{def:g11:functional-groups:halogenoalkane}{haloalkana}; \omterm{def:g11:functional-groups:halogenoalkane}{alkena}.
\end{solution}

\begin{solution}{exo:g11:functional-groups:4}
Propanal \ce{CH3-CH2-CHO}, propanon \ce{CH3-CO-CH3}. Propanal adalah
\omterm{def:g11:functional-groups:carbonyl}{aldehida}: karbon \ce{C=O}-nya berada di ujung rantai, terikat pada
sebuah \omterm{def:g7:atoms-and-molecules:atom}{atom} hidrogen; pada propanon karbon itu berada di antara dua \omterm{def:g7:atoms-and-molecules:atom}{atom}
karbon.
\end{solution}

\begin{solution}{exo:g11:functional-groups:5}
Primer (karbon yang mengikat \ce{-OH} memiliki satu tetangga karbon);
sekunder (dua); tersier (tiga).
\end{solution}

\begin{solution}{exo:g11:functional-groups:6}
\setchemfig{atom sep=1.5em}
Pentan-3-ol \chemfig{-[:30]-[:-30](-[:-90]OH)-[:30]-[:-30]};
asam 2-metilbutanoat \chemfig{-[:30]-[:-30](-[:-90])-[:30](=[:90]O)-[:-30]OH};
heksan-2-on \chemfig{-[:30](=[:90]O)-[:-30]-[:30]-[:-30]-[:30]};
propil etanoat \chemfig{-[:30](=[:90]O)-[:-30]O-[:30]-[:-30]-[:30]}.
\end{solution}

\begin{solution}{exo:g11:functional-groups:7}
\ce{C3H6O}. Ya, keduanya \omterm{def:g11:organic-skeletons:isomer}{isomer}, tetapi dengan gugus yang berbeda:
\omterm{def:g11:functional-groups:carbonyl}{aldehida} dan \omterm{def:g11:functional-groups:carbonyl}{keton} (keduanya gugus karbonil, pada tempat yang berbeda).
\end{solution}

\begin{solution}{exo:g11:functional-groups:8}
\ce{HCOO-CH2-CH3}, etil metanoat; \ce{CH3-COO-CH2-CH2-CH3}, propil
etanoat.
\end{solution}

\begin{solution}{exo:g11:functional-groups:9}
Metil etanoat (asam etanoat, metanol); etil propanoat (asam propanoat,
etanol); metil metanoat (asam metanoat, metanol).
\end{solution}

\begin{solution}{exo:g11:functional-groups:10}
\setchemfig{atom sep=1.5em}
\begin{center}
\small
\begin{tabular}{cccc}
\chemfig{-[:30]-[:-30]-[:30]-[:-30]OH} &
\chemfig{-[:30](-[:90]OH)-[:-30]-[:30]} &
\chemfig{-[:30](-[:90])-[:-30]-[:30]OH} &
\chemfig{-[:0](-[:90]OH)(-[:-90])-[:0]} \\
butan-1-ol & butan-2-ol &
2-metilpropan-1-ol & 2-metilpropan-2-ol \\
primer & sekunder & primer & tersier
\end{tabular}
\end{center}
\end{solution}

\begin{solution}{exo:g11:functional-groups:11}
Sebuah \omterm{def:g11:functional-groups:amine}{amida} (etanamida). \emph{-ena}. \omterm{def:g11:functional-groups:carboxylic-acid}{Asam karboksilat} (\ce{C=O} dan
\ce{-OH} pada karbon yang sama), \omterm{def:g11:functional-groups:carboxylic-acid}{ester} (\ce{C=O} dan \ce{-O-}), dan \omterm{def:g11:functional-groups:amine}{amida}
(\ce{C=O} dan \ce{N}): tiga keluarga.
\end{solution}

\begin{solution}{exo:g11:functional-groups:12}
Aspirin: gugus karboksil (\omterm{def:g11:functional-groups:carboxylic-acid}{asam karboksilat}) dan gugus \omterm{def:g11:functional-groups:carboxylic-acid}{ester}.
Parasetamol: gugus \ce{-OH} fenol pada cincin dan gugus \omterm{def:g11:functional-groups:amine}{amida}
\ce{NH-CO}.
\end{solution}

\begin{solution}{exo:g11:functional-groups:13}
\ce{C2H6O}. Hanya etanol, dengan \ce{O-H}-nya, yang membentuk \omterm{def:g11:polarity-and-cohesion:hydrogen-bond}{ikatan
hidrogen} di antara \omterm{def:g7:atoms-and-molecules:molecule}{molekul-molekulnya} (metoksimetana tidak memiliki
hidrogen pada oksigennya). Etanol mendidih lebih tinggi, pada
\qty{78}{\celsius}; metoksimetana pada \qty{-25}{\celsius}.
\end{solution}

\begin{solution}{exo:g11:functional-groups:14}
Gugus hidroksil dan gugus karboksil. Rantainya memiliki tiga karbon,
dengan karbon asam sebagai karbon 1 dan \ce{-OH} pada karbon 2: asam
2-hidroksipropanoat. Glisina: asam 2-aminoetanoat (sering ditulis asam
aminoetanoat).
\end{solution}

\begin{solution}{exo:g11:functional-groups:15}
Etanol \ce{CH3-CH2-OH}, \ce{C2H6O}, \omterm{def:g11:functional-groups:alcohol}{alkohol}; etanal \ce{CH3-CHO},
\ce{C2H4O}, \omterm{def:g11:functional-groups:carbonyl}{aldehida}; asam etanoat \ce{CH3-COOH}, \ce{C2H4O2}, \omterm{def:g11:functional-groups:carboxylic-acid}{asam
karboksilat}. Dari etanol ke etanal, dua \omterm{def:g7:atoms-and-molecules:atom}{atom} hidrogen hilang; dari
etanal ke asam etanoat, satu \omterm{def:g7:atoms-and-molecules:atom}{atom} oksigen bertambah.
\end{solution}

\begin{solution}{pb:g11:functional-groups:1}
\textbf{1.} A: gugus \omterm{def:g11:functional-groups:carboxylic-acid}{ester}; B: hidroksil; C: karboksil; D: karbonil di
ujung rantai; E: gugus \omterm{def:g11:functional-groups:carboxylic-acid}{ester}.

\textbf{2.} A \omterm{def:g11:functional-groups:carboxylic-acid}{ester}, B \omterm{def:g11:functional-groups:alcohol}{alkohol}, C \omterm{def:g11:functional-groups:carboxylic-acid}{asam karboksilat}, D \omterm{def:g11:functional-groups:carbonyl}{aldehida}, E \omterm{def:g11:functional-groups:carboxylic-acid}{ester}.

\textbf{3.} A (pisang) dan E (nanas): bau buah.

\textbf{4.} Primer: karbon yang mengikat \ce{-OH} terikat pada satu
karbon.

\textbf{5.} \ce{C6H12O}; misalnya heksan-2-on,
\ce{CH3-CO-CH2-CH2-CH2-CH3}.

\textbf{6.} Rantai yang memuat \ce{C-OH} memiliki empat karbon, dinomori
dari \ce{-OH}, dengan metil pada karbon 3: 3-metilbutan-1-ol.

\textbf{7.} C: asam etanoat; D: heksanal.

\textbf{8.} Etil butanoat, kerabat asam butanoat dan etanol.

\textbf{9.} \emph{Etanoat}: bagian asam \ce{CH3-CO}, dari asam etanoat;
\emph{3-metilbutil}: bagian \omterm{def:g11:functional-groups:alcohol}{alkohol}, dari 3-metilbutan-1-ol.

\textbf{10.} C (asam etanoat) dan B (3-metilbutan-1-ol).

\textbf{11.} A \ce{C7H14O2}; B \ce{C5H12O}; C \ce{C2H4O2}.

\textbf{12.} $\ce{C2H4O2} + \ce{C5H12O}$ mengandung \ce{C7H16O3}, satu
\ce{H2O} lebih banyak daripada \ce{C7H14O2}: air.

\textbf{13.} \ce{C2H4O2 + C5H12O -> C7H14O2 + H2O}: C 7 = 7, H 16 = 16,
O 3 = 3.

\textbf{14.} $M(\text{B}) = 5 \times 12.0 + 12 \times 1.0 + 16.0 =
\qty{88.0}{g/mol}$; $M(\text{C}) = \qty{60.0}{g/mol}$.

\textbf{15.} $60.0 + 88.0 - 18.0 = \qty{130.0}{g/mol}$.

\textbf{16.} $7 \times 12.0 + 14 \times 1.0 + 2 \times 16.0 =
\qty{130.0}{g/mol}$: sama.

\textbf{17.} E, \ce{C6H12O2}: $6 \times 12.0 + 12 \times 1.0 + 2 \times
16.0 = \qty{116.0}{g/mol}$.

\textbf{18.} \ce{C7H14O2}: rumusnya sama dengan A, jadi \omterm{def:g11:organic-skeletons:isomer}{isomer}, dengan
\omterm{def:g10:the-mole:molar-mass}{massa molar} yang sama. \omterm{def:g10:the-mole:molar-mass}{Massa molar} saja tidak dapat membedakan
\omterm{def:g11:organic-skeletons:isomer}{isomer-isomer}.

\textbf{19.} A dan asam heptanoat memiliki \omterm{def:g7:atoms-and-molecules:atom}{atom-atom} yang sama tetapi
gugus yang berbeda, \omterm{def:g11:functional-groups:carboxylic-acid}{ester} dan \omterm{def:g11:functional-groups:carboxylic-acid}{asam karboksilat}: yang satu berbau pisang,
yang lain tengik. Gugusnya, dan letaknya, yang menentukan.

\textbf{20.} \qty{130.0}{g/mol}.
\end{solution}
